\section{Sovereign Governance}

\label{sec:governance}
% ============================================================

\subsection{Governance Instantiation}

Each chain $\Chain_i$ that connects to Teleport deploys an OmnichainRouter instance $R_i$ with a designated governor $\Gov_i$. The governor is typically:

\begin{itemize}[nosep]
  \item \textbf{Lux C-Chain}: LUX token DAO with 48-hour timelock. Fees flow to xLUX holders.
  \item \textbf{Zoo EVM}: ZOO token governance. Fees flow to ZOO stakers.
  \item \textbf{External L1/L2}: The chain's own governance mechanism, chosen at deployment.
\end{itemize}

\noindent
The key insight is that $\Gov_i$ controls \emph{only} the economic parameters of $R_i$. It has no authority over $R_j$ for $j \neq i$, and no authority over the MPC signer set or immutable protocol logic.

\subsection{Governor Capability Matrix}

\begin{table}[h]
\centering
\small
\begin{tabular}{@{}p{3.8cm}cc@{}}
\toprule
\textbf{Action} & \textbf{CAN} & \textbf{CANNOT} \\
\midrule
Set bridge fee (0--100 bps) & $\checkmark$ & \\
Set daily mint limits & $\checkmark$ & \\
Register new bridge tokens & $\checkmark$ & \\
Enable Shariah mode & $\checkmark$ & \\
Set yield strategy routing & $\checkmark$ & \\
Change fee recipient (48h lock) & $\checkmark$ & \\
Transfer governance (2-step) & $\checkmark$ & \\
\midrule
Mint bridge tokens & & $\times$ \\
Modify nonce bitmaps & & $\times$ \\
Set fee $> 100$ bps & & $\times$ \\
Bypass signature verification & & $\times$ \\
Upgrade contract bytecode & & $\times$ \\
Override signer rotation & & $\times$ \\
Access other chains' $R_j$ & & $\times$ \\
\bottomrule
\end{tabular}
\caption{Governor capability matrix. The ``CANNOT'' column represents hard-coded invariants enforced at the bytecode level, not merely policy.}
\label{tab:capability}
\end{table}

\subsection{Fee Configuration}

The governor sets the bridge fee rate $f_b$ and stakeholder share $f_s$:

\begin{equation}
  0 \leq f_b \leq 100 \text{ bps}, \quad 0 \leq f_s \leq 10{,}000 \text{ bps}
\end{equation}

The hard cap $f_b \leq 100$ bps (1\%) is enforced at the bytecode level:

\begin{lstlisting}[language=Solidity,caption=Fee cap enforcement]
function setBridgeFee(uint256 newFeeBps)
    external onlyGovernor
{
    require(
        newFeeBps <= MAX_FEE_BPS,
        "fee exceeds hard cap"
    );
    // MAX_FEE_BPS = 100, immutable
    bridgeFeeBps = newFeeBps;
    emit BridgeFeeUpdated(newFeeBps);
}
\end{lstlisting}

\noindent
The constant \texttt{MAX\_FEE\_BPS = 100} is defined at compile time and cannot be modified post-deployment, even by the governor.

\subsection{Fee Distribution}

For a bridge operation of amount $a$ on chain $\Chain_i$ with fee rate $f_b^{(i)}$ and stakeholder share $f_s^{(i)}$:

\begin{align}
  \text{fee} &= a \cdot \frac{f_b^{(i)}}{10{,}000} \\
  \text{toStakeholders} &= \text{fee} \cdot \frac{f_s^{(i)}}{10{,}000} \\
  \text{toTreasury} &= \text{fee} - \text{toStakeholders}
\end{align}

Each chain instance independently determines where fees flow:
\begin{itemize}[nosep]
  \item On Lux C-Chain: stakeholder vault $=$ xLUX (LiquidLUX).
  \item On Zoo EVM: stakeholder vault $=$ ZOO staking contract.
  \item On white-label deployments: stakeholder vault $=$ deployer-designated address.
\end{itemize}

\subsection{Signer Rotation Protection}

MPC signer rotation is the highest-severity operation. The rotation delay is a configurable operational parameter (default 24h, governor-settable $\in [0, 7\text{d}]$), not a security invariant---the MPC threshold \emph{is} the trust model. A compromised governor cannot initiate rotation; only the current $t$-of-$n$ MPC set can.

\begin{enumerate}[nosep]
  \item Current $t$-of-$n$ MPC set signs the rotation proposal with the new MPC group address.
  \item Configurable delay begins (default 24h). Pending signers are published on-chain.
  \item During the delay, users may exit by burning bridge tokens for underlying assets.
  \item After the delay expires, any address may execute the rotation.
  \item Current signers may cancel during the delay period.
\end{enumerate}

\begin{property}[Exit Window]
\label{prop:exit}
For any signer rotation event at time $t_0$ with rotation delay $\delta$ (default 24h, max 7 days), every bridge token holder has the interval $[t_0, t_0 + \delta)$ to exit the system by burning bridge tokens and receiving the underlying asset at the current share price.
\end{property}

% ============================================================
